\documentclass[10pt,a4paper]{amsart}
\usepackage[margin=19mm]{geometry}
\usepackage{amsmath,amssymb,mathtools}
\usepackage[scr=rsfs,cal=euler]{mathalfa}
\usepackage{xfrac}
\usepackage[unicode,hidelinks,bookmarks=false]{hyperref}
\newcommand{\ie}{i.e.\ }
\newcommand{\cf}{cf.\ }
\newcommand{\resp}{respectively\ }
\newcommand{\Subsection}{\subsection}
\providecommand{\nobreakdash}{\nobreak\hbox{-}}
\setlength{\emergencystretch}{2em}
\multlinegap0pt
\usepackage{bm}
\usepackage{bbm}
\usepackage{dsfont}
\usepackage{xspace}
\usepackage{cancel}
\usepackage{mathtools}
\usepackage{amsthm}
\makeatletter
\@ifundefined{theorem}{\newtheorem{theorem}{Theorem}}{}
\@ifundefined{proposition}{\newtheorem{proposition}[theorem]{Proposition}}{}
\makeatother
\newcommand{\RtoL}{\text{RevtoL}\xspace} % algorithm Revoke to the left
\title[Online Interval Selection on a Simple Chain]{Online Interval Selection on a Simple Chain}
\author{Yaqiao Li, Ali Mohammad Lavasani, Denis Pankratov}
\date{Source: arXiv:2608.10376v1 (CC BY 4.0)}
\begin{document}
\maketitle

\noindent\emph{Source and licence.} Online Interval Selection on a Simple Chain. Version arXiv:2608.10376v1 (CC BY 4.0), distributed under the Creative Commons Attribution 4.0 International licence, as recorded in the arXiv licence metadata for this exact version and shown on the abstract page https://arxiv.org/abs/2608.10376v1. Full rights notice: licenses/arxiv-2608.10376-CC-BY-4.0.txt.\par\smallskip

\noindent\textit{Editorial scope.} Counted unit: the Proposition whose source label is \texttt{thm:formula\_dn\_difference} in the article (source0.tex lines 206--209, source label \texttt{thm:formula\_dn\_difference}), delivered together with the article's own complete original proof of it (source0.tex lines 211--267, one \texttt{proof} environment of 57 lines). No mathematics is written, repaired or re-derived: statement and proof are the article's own text line for line. Required context retained uncounted: none. Every definition, display and label of those lines is retained unchanged. Omitted from this package and not redistributed: 1--205, 210--210, 268--481. Nothing inside a retained excerpt is omitted or altered: every retained body is delivered exactly as the article's lines are written, so no omission receipt of any kind accompanies this package. Citations: the retained excerpts carry no citation command at all, so no bibliography entry of the article is transcribed and no reference is redistributed. Reference adaptation: none was needed and none was made; every reference inside the delivered unit resolves inside it, so no unresolved reference or citation is produced. Printed slips kept literal: no printed word, symbol, hypothesis or number of the counted statement or of its proof was altered. Layout and class: the article's own class is \texttt{article} with packages including geometry, amssymb, amsmath, amsthm, xspace, tikz, hyperref; the wrapper is the lane's pinned \texttt{amsart} editorial wrapper at 10pt on A4, which restates the theorem-like environments the retained text needs and adds the article's own macro definitions that the retained text needs (\texttt{\textbackslash RtoL}), with extra wrapper packages (bm, bbm, dsfont, xspace, cancel, mathtools). The delivered numbering is the unit's own. Assets: no retained span contains a figure environment, a graphics-inclusion command, a PDF-page inclusion, a table or a TikZ picture; no image file is copied into this package, the package declares no asset dependency and no image file is redistributed with it. Licence: version arXiv:2608.10376v1 of this article is distributed under the Creative Commons Attribution 4.0 International licence, as recorded in the arXiv licence metadata for this exact version and shown on the abstract page https://arxiv.org/abs/2608.10376v1. A full rights notice is delivered at licenses/arxiv-2608.10376-CC-BY-4.0.txt. Characters: every retained excerpt is pure ASCII, so the delivered engine, XeLaTeX with its default Latin Modern text font, reports no missing character.
\begin{proposition}   \label{thm:formula_dn_difference}
    For $n\ge 4$,
        $d_{n,n} - d_{n,n-1} = 2d_{n-1,n-2} - d_{n-1,n-1} - d_{n-1,n-3}$.
\end{proposition}

\begin{proof}
    We give a formula for $d_n$. For 
        $1 \le i \le n$, 
    define
        \begin{equation}    \label{eq:def_pi}
	\begin{split}
            p_i &= \#\{\sigma \in S_n: \text{ interval } n \text{ is  chosen by \RtoL,}  \\
			&\text{ while interval } i \text{ is the last interval of the input} \}.
	\end{split}
        \end{equation}
    Then,
        $d_n = p_1 + \ldots + p_n$.
    Below we give formulas for computing $p_i$.

 	(i) $p_n = (n-1)! - d_{n-1}$.             This is because interval $n$ is chosen iff interval $n-1$ is \emph{not} chosen among intervals
                $1, \ldots, (n-1)$.

       (ii)  $p_{n-1} = (n-1) d_{n-2} = d_{n-1,n-2}$.       This is because in this case interval $n$ is chosen if and only if interval $n-2$ is chosen among intervals 
                $1, \ldots, (n-2)$, 
            and we have the multiplicative factor $n-1$ because interval $n$ can be arbitrarily placed in any of the $n-1$ positions.

        (iii)  for $2 \le i \le n-2$, 
            $p_i = (n-1)\cdots(n-i+1)d_{n-i} = d_{n-1,n-i}$.
            
            Consider the two blocks of intervals: intervals 
                $1, \ldots, (i-1)$, 
            and intervals 
                $(i+1), \ldots, n$. 
            Since interval $i$ is not present when these $n-1$ intervals appear, we know that these two blocks of intervals do not interfere with each other. Hence, interval $n$ is chosen if and only if it is chosen among intervals 
                $(i+1), \ldots, n$. 
            But this is equivalent to interval $n-i$ is chosen among intervals 
                $1, \ldots, (n-i)$. 
            Since the other block of intervals 
                $1, \ldots, (i-1)$ 
            can be arbitrarily placed in $n-1$ positions, we have the multiplicative factor 
                $(n-1)\cdots(n-i+1)$.

        (iv)  $p_1 = d_{n-1}$.      This is because in this case interval $n$ is chosen if and only if interval $n$ is chosen among intervals 
                $2, \ldots, n$, 
            which is equivalent to interval $n-1$ is chosen among intervals 
                $1, \ldots, (n-1)$.

Hence, 
        $d_n
        = (n-1)! + 2d_{n-1,n-2} + \sum_{j=2}^{n-3} d_{n-1,j}$.
    With this, we have %(apply Lemma\ref{lem:one_step_reduction} whenever necessary)
\begin{align*}
        &d_{n,n} - d_{n,n-1}= d_{n,n} - n d_{n-1,n-1} = d_{n,n} - (n-1) d_{n-1,n-1} - d_{n-1,n-1} \\
        &= \left( (n-1)! + 2d_{n-1,n-2} + \sum_{j=2}^{n-3} d_{n-1,j} \right) \\
        &\phantom{==} - (n-1) \left( (n-2)! + 2d_{n-2,n-3} + \sum_{j=2}^{n-4} d_{n-2,j} \right) 
        - d_{n-1,n-1}   \\
        &= \left( (n-1)! + 2d_{n-1,n-2} + \sum_{j=2}^{n-3} d_{n-1,j} \right) \\
        &\phantom{==} - \left( (n-1)! + 2d_{n-1,n-3} + \sum_{j=2}^{n-4} d_{n-1,j} \right)
        - d_{n-1,n-1}   \\
        &= 2d_{n-1,n-2} - d_{n-1,n-1} - d_{n-1,n-3}.	\qedhere
\end{align*}
\end{proof}

\end{document}
